\documentclass{article}
\usepackage{amsmath,hyperref}
\title{Rest-to-rest formation recovery when an anchor change is learned late}
\date{}
\begin{document}
\maketitle

\paragraph{Physical and information laws.}
A stopped vehicle splits an indoor inspection corridor into left and right crews. Each has one moving terminal robot; their radio link can bridge the gap when their final separation is at most the licensed range. Inward travel on each side has a physical cap from clearance of the stopped body and continued coverage of that side's assigned inspection anchor. The anchor need not remain at the old formation position. In this model each robot moves monotonically inward, has bounded forward speed $V$, positive maximum braking $b$ and nonnegative acceleration limit $a$, and must be stationary at deadline $T$. The clearance and anchor limits apply throughout, not just at $T$. Radio connectivity may be broken before the deadline but must be restored by it. No curvature, turning, unmodeled motor lag, jerk, shared sample-clock protocol or radio power model is inferred from this deliberately narrow envelope. A missing or late observation cannot be supplied by the proof-side configuration.

\paragraph{Known local cap.}
Let $v_0$ be the initial inward speed, with $0\le v_0\le V$, and let $M$ be the smaller clearance and anchor travel cap. If $T<v_0/b$, rest is impossible even without a physical cap. If $M<v_0^2/(2b)$, physical clearance or anchor coverage cannot be preserved even with maximum braking; this is an initial state outside the viable region. Otherwise exact feasible travel is the interval
\[
\left[\frac{v_0^2}{2b},
      \min\left(M,\int_0^T\min\{V,v_0+at,b(T-t)\}\,dt\right)\right].
\]
Every admissible velocity is bounded below by $v_-(t)=\max(0,v_0-bt)$ and above by $v_+(t)=\min(V,v_0+at,b(T-t))$, using the respective initial, terminal and actuation bounds. These piecewise-affine profiles themselves have the required endpoints and slopes. Their convex interpolation gives every travel between their integrals, and its nonnegative velocity ensures position never exceeds its final value. This is a continuous-time witness, not a terminal displacement ignoring acceleration or an implicit interpolation of integer samples.

Knowing each side's own cap is enough for the local maximum-travel policies to be optimal for the joint range condition: they add their respective maximum permitted inward travels without needing the other's reserve. Their sum must meet the separation deficit. Before treating this as restoration, the envelope also checks both sides' initial braking feasibility. This simple decomposition breaks if the caps are not known to the actors at decision times or if an action changes the other's legal region.

\paragraph{One delayed observation.}
For an exact two-world discriminator, take initial speed zero, unit acceleration and unit braking, and a speed cap large enough not to bind. A local change in anchoring duties implies the terminal robot's total permitted travel is either $M$ or a larger value, but it learns which only at the usable board receipt $\tau$ (including reaction time). The same prefix must be safe and viable in both worlds, and the cutoff and rest time are fixed in this family. A missing distinguishing result keeps the tight cap in force. The board's abstract observation is truthful by the declared source law; a new device or fault that changes what its bit warrants must be revalidated. The comparison uses the realized authenticated result at its actual cut, not the globally true mode ahead of time.

Write $R=T-\tau$, and let $z$ and $x$ be the speed and distance traveled by the cut. For each $0\le z\le\min(\tau,R,\sqrt M)$ the possible prefix distance lies between $z^2/2$ and
\[
G(\tau,z)=(\tau^2+2\tau z-z^2)/4.
\]
These come from late acceleration and from the integral of the maximal prefix velocity $\min(t,z+\tau-t)$, respectively. Linear interpolation of those velocities realizes intermediate distances. Even at the cut the tight alternative must have room to brake; hence $x+z^2/2\le M$. This adds exactly the upper constraint $x\le M-z^2/2$ and introduces the bound $z\le\sqrt M$. Nondecreasing position and an executable tight-world brake ensure this prefix was also safe along the way. For fixed $z$, take the largest feasible $x=\min(G(\tau,z),M-z^2/2)$; it weakly increases both worlds' final reachable distances. The analogous maximal suffix distance, with endpoint speeds $z,0$, is $G(R,z)$. Thus the greatest unclipped total at the deadline is exactly
\[
 B(\tau)=\max_{0\le z\le\min(\tau,R,\sqrt M)}
       \left[\min\{G(\tau,z),M-z^2/2\}+G(R,z)\right].
\]
The reachable maximum in each eventual world is the lesser of its cap and $B(\tau)$. The same optimal prefix achieves all these maxima simultaneously; the suffix then uses the informed legal cap. This assertion holds here because both worlds have the same actuators and target horizon, and differ only in caps. It would not follow for arbitrary different actuation laws or safety predicates.

The switch between the two pieces is at $z=2\sqrt M-\tau$. The free piece is nondecreasing because $z\le T/2$. The braking-limited piece is concave with optimum at $z=R/3$, clipped to its feasible segment. The replay evaluates these candidates exactly for rational $\sqrt M$. This is an analytic maximization, not a claim based on the finite sample. For $T=4$, $M=1$, one gets $B(\tau)=4$ for $0\le\tau\le1$, by following the unconstrained optimal prefix and retaining enough room to stop. For $1\le\tau\le4$, completing the square bounds the expression by $1+(4-\tau)^2/3$, attained by the braking-limited prefix at $z=(4-\tau)/3$. At $\tau=2$ the bound is $7/3$, with $x=7/9$ and $z=2/3$. Lean checks the algebraic upper bound, attaining cut state, exact root cutoff and endpoint braking inequality, not the full continuous integration argument.

\paragraph{A matched locally impossible join.}
Use an initial radio separation deficit of three distance units with enough geometric clearance for the displacements below. In the first source history the left cap is one and the right crew can contribute two; in the second the left cap exceeds four and the right crew must stay with its own anchor. The right crew knows its own current anchor locally. The left crew's local history is identical until its board receipt; it cannot assume the right crew will provide spare travel in the loose world. This paired assignment is one given source law, not independent resampling of both caps. Both histories have perfectly safe joining trajectories when known before motion. A common local policy with a timely board result can join in both exactly when $\tau \le 4-\sqrt6$. A result at tick one meets both duties; a result at tick two still permits collision-free braking and the tight-world join but cannot close the loose-world gap by $T=4$. An erased result cannot guarantee both histories. In the failure branch the safe containment policy still brakes; it is not reported as meeting the radio deadline. Independently combining the tight cap with the immobile right crew creates a different law with physical impossibility, so that ablation is not a stronger matched baseline.

\paragraph{Comparators and evidence.}
\href{https://www.sciencedirect.com/science/article/pii/S0016003223003046}{Maeda, Endo and Matsuno} already address decentralized navigation with the failed unit remaining in the environment. \href{https://en.endo-lab.jp/publication/detail/12147}{Endo et al. 2026} describe distributed model predictive control with limited observation, collision avoidance and connectivity recovery. We verified the authors' and IEEE abstracts, not their detailed constraints or code, and do not claim an implementation comparison with them. \href{https://arxiv.org/abs/2210.01354}{Breeden and Panagou} address simultaneous barrier feasibility with bounded inputs. The exact advance in this particular fixture is to isolate continuous stopping and a late-information service bound, not to supersede barrier methods or MPC.

The matched trajectory-LP reference has the same dynamics and information cut, joins the prefix controls before the receipt and branches only afterward. On an independent uniform segmentation it is a lower bound on the exact continuous optimum and matched it to within the declared tolerance on the tested cases. The known-cap LP has analytic breakpoints supplied, so the current replay verifies the construction and numerical optimizer but is not a separate discovery of those breakpoints. The replays passed 520 known-cap numerical comparisons with exact constructive witnesses, 2,600 rational two-world cases and 90 independent-grid LP checks. The one-dimensional synthetic result is a partial R6 comparison, not full obligation acceptance or field assurance. Larger teams, unreliable visits after observation, and an independently evaluated source device remain separate.
\end{document}
